% !TEX root = ../main.tex
\section{Introduction}
\label{sec:introduction}
Changing an answer can indicate correction or unjustified concession;
maintaining it can indicate justified resistance or persistence in error.
Neither behaviour is interpretable without the task and evidence. A jury
adds a further complication: it can endorse a defence whose conclusion
disagrees with the source reference.

We reuse benchmark work and add a dispute after completion. The model answers
and states its confidence before Veracity is revealed. We freeze that output,
branch it into four challenges, and separately measure revision, forecasts of
correctness and forecasts of jury acceptance. The jury judges the original;
a corrected replacement does not retroactively repair it.

Track~1 of the AI x Epistemics Research Sprint~\cite{apart2026} motivates three
questions: do prospective probabilities predict source agreement; does pressure
erode matching answers while audit repairs mismatches; and can an evaluator
leak its target? We test the last question offline, without attributing
exploitation to participants. Track~2 motivates auditable verification.

Five model families receive every selected case across five benchmarks and
four occupational domains. We retain literal and audited grades because units,
rounding and fallible references affect apparent errors. Four selected exchanges
illustrate how these measures can diverge. The sampled domains do not establish
representative industry performance. The preceding pilot remains a separate cohort.
